\subsection{Embedded Runge--Kutta formulas and adaptive step size}
\label{sec:RKAdaptive}

Hairer I, Section~II.4, printed pp.~165--169, Eqs.~(4.6)--(4.14)
\cite{HNW1993}. The calculation proceeds in this order:
\[
 \text{shared stages}\ \longrightarrow\ \text{two endpoint values}
 \ \longrightarrow\ \text{error vector},
\]
\[
 \text{error vector}\ \longrightarrow\ \text{scaled scalar error}
 \ \longrightarrow\ \text{accept or retry and choose the next step}.
\]
We use boldface for state and error vectors. Let
$\mathbf y'=\mathbf f(x,\mathbf y)$, $\mathbf y\in U\subseteq\mathbb R^d$,
with $d\geq1$. Assume the smoothness needed for the stated local orders on a
neighborhood containing the stages. Initially take $h>0$. All error comparisons
refer to a single attempted step from the \emph{same} current state
$(x_n,\mathbf y_n)$, with exact local reference
\begin{equation}\label{Eq:AdaptLocalReference}
 \mathbf Y=\boldsymbol\Psi_h(x_n,\mathbf y_n).
\end{equation}
Even if $\mathbf y_n$ contains accumulated error, $\mathbf Y$ starts from that
computed state, not from the original trajectory's exact state at $x_n$.

\subsubsection{One set of stages, two sets of output weights}
For an explicit $s$-stage pair, compute
\begin{equation}\label{Eq:AdaptSharedStages}
\begin{aligned}
 \mathbf k_1&=\mathbf f(x_n,\mathbf y_n),\\
 \mathbf k_i&=\mathbf f\!\left(x_n+c_i h,
       \mathbf y_n+h\sum_{j<i}a_{ij}\mathbf k_j\right),
       \quad i=2,\ldots,s,\\
 c_i&=\sum_{j<i}a_{ij}.
\end{aligned}
\end{equation}
Use these same derivative values in two updates:
\begin{equation}\label{Eq:AdaptTwoOutputs}
\begin{aligned}
 \mathbf y_{n+1}^{\mathrm{high}}
 &=\mathbf y_n+h\sum_i b_i\mathbf k_i,\\
 \mathbf y_{n+1}^{\mathrm{low}}
 &=\mathbf y_n+h\sum_i\widetilde b_i\mathbf k_i.
\end{aligned}
\end{equation}
Subtract the second update from the first, displaying the common term:
\begin{equation}\label{Eq:AdaptDifference}
\begin{aligned}
 \boldsymbol\delta
 &:=\mathbf y_{n+1}^{\mathrm{high}}-\mathbf y_{n+1}^{\mathrm{low}}\\
 &=\left[\mathbf y_n+h\sum_i b_i\mathbf k_i\right]
   -\left[\mathbf y_n+h\sum_i\widetilde b_i\mathbf k_i\right]\\
 &=h\sum_i(b_i-\widetilde b_i)\mathbf k_i.
\end{aligned}
\end{equation}
No new right-hand-side evaluation is needed to form this difference once the
pair's stages are available.

Let the lower formula have order $q$ and the higher formula order $p>q$.
At fixed starting data, write their signed local errors as
\[
 \mathbf Y-\mathbf y_{n+1}^{\mathrm{low}}
   =\mathbf L h^{q+1}+O(h^{q+2}),\qquad
 \mathbf Y-\mathbf y_{n+1}^{\mathrm{high}}=O(h^{p+1}).
\]
Then
\begin{equation}\label{Eq:AdaptEstimatedOrder}
\begin{aligned}
 \boldsymbol\delta
 &=(\mathbf Y-\mathbf y_{n+1}^{\mathrm{low}})
    -(\mathbf Y-\mathbf y_{n+1}^{\mathrm{high}})\\
 &=\mathbf L h^{q+1}+O(h^{q+2})-O(h^{p+1})\\
 &=\mathbf L h^{q+1}+O(h^{q+2}),
\end{aligned}
\end{equation}
since $p+1\geq q+2$. Thus $\boldsymbol\delta$ estimates the lower-order
local error. The solver can still keep the higher-order value; HNW calls
this choice \emph{local extrapolation}. It does not mean adding a Richardson
step-doubling correction.

\subsubsection{A two-stage example: Euler embedded in Heun}
Use
\[
 \mathbf k_1=\mathbf f(x_n,\mathbf y_n),\qquad
 \mathbf k_2=\mathbf f(x_n+h,\mathbf y_n+h\mathbf k_1),
\]
with output weights $b=(1/2,1/2)$ and $\widetilde b=(1,0)$:
\begin{equation}\label{Eq:AdaptHeunEuler}
\begin{aligned}
 \mathbf y_{n+1}^{\mathrm{high}}
   &=\mathbf y_n+\tfrac h2(\mathbf k_1+\mathbf k_2),\\
 \mathbf y_{n+1}^{\mathrm{low}}&=\mathbf y_n+h\mathbf k_1,\\
 \boldsymbol\delta
   &=\tfrac h2(\mathbf k_1+\mathbf k_2)-h\mathbf k_1
     =\tfrac h2(\mathbf k_2-\mathbf k_1).
\end{aligned}
\end{equation}
For $y'=y$ starting at $y_n=1$,
\[
\begin{aligned}
 k_1&=1,\qquad k_2=1+h,\\
 y_{n+1}^{\mathrm{low}}&=1+h,\\
 y_{n+1}^{\mathrm{high}}&=1+h+\tfrac12h^2,\\
 \delta&=\tfrac h2[(1+h)-1]=\tfrac12h^2.
\end{aligned}
\]
The exact local value is $Y=e^h$, so
\[
\begin{aligned}
 Y-y_{n+1}^{\mathrm{low}}
 &=\tfrac12h^2+\tfrac16h^3+O(h^4)=\delta+O(h^3),\\
 Y-y_{n+1}^{\mathrm{high}}&=\tfrac16h^3+O(h^4).
\end{aligned}
\]
The estimate starts at $h^2$, although the retained Heun value has local
error starting at $h^3$.

\subsubsection{DOPRI5: what the difference estimates}
For the Dormand--Prince $5(4)$ pair, $p=5$, $q=4$, $s=7$; see Hairer I,
Section~II.5, printed pp.~178--179, Table~5.2 \cite{HNW1993}.
Its fifth-order weights are
\[
 b=\left(\frac{35}{384},0,\frac{500}{1113},\frac{125}{192},
 -\frac{2187}{6784},\frac{11}{84},0\right).
\]
Define the error weights $\varepsilon_i=b_i-\widetilde b_i$:
\[
 \varepsilon=\left(\frac{71}{57600},0,-\frac{71}{16695},
 \frac{71}{1920},-\frac{17253}{339200},\frac{22}{525},-\frac1{40}\right).
\]
Equation~\ref{Eq:AdaptDifference} then becomes
\begin{equation}\label{Eq:AdaptDopriDifference}
\begin{aligned}
 \boldsymbol\delta
 &=h\left(\frac{71}{57600}\mathbf k_1
       -\frac{71}{16695}\mathbf k_3+\frac{71}{1920}\mathbf k_4\right.\\
 &\hspace{2.2em}\left.
       -\frac{17253}{339200}\mathbf k_5+\frac{22}{525}\mathbf k_6
       -\frac1{40}\mathbf k_7\right),\\
 \boldsymbol\delta&=\mathbf L h^5+O(h^6),\qquad
 \mathbf Y-\mathbf y_{n+1}^{\mathrm{high}}=O(h^6).
\end{aligned}
\end{equation}
The seventh row satisfies $a_{7j}=b_j$ and $c_7=1$, so
\[
 \mathbf k_7=\mathbf f(x_n+h,\mathbf y_{n+1}^{\mathrm{high}}).
\]
After acceptance it can be reused as the next step's $\mathbf k_1$
(\emph{first same as last}, FSAL). The first attempt requires seven stage evaluations, apart from any
starting-step probe; subsequent attempts reuse the derivative at their
starting state and need six new evaluations. After a rejection, the starting state and its derivative
remain unchanged; the rejected endpoint's derivative is not the next
starting derivative.

\subsubsection{Scaling each error component by its tolerance}
A state can combine positions, velocities, and other quantities with different
units and magnitudes. Set, for each component $r$,
\begin{equation}\label{Eq:AdaptToleranceScale}
 \mathrm{sc}_r=\mathrm{Atol}_r
 +\mathrm{Rtol}_r\max\!\left(|y_n^r|,
                      |(y_{n+1}^{\mathrm{high}})^r|\right),
 \qquad z_r=\frac{\delta^r}{\mathrm{sc}_r}.
\end{equation}
Here $\mathrm{Atol}_r$ has the units of $y^r$, while $\mathrm{Rtol}_r$ and
$z_r$ are dimensionless. Assume finite
$\mathrm{Atol}_r>0$ and finite $\mathrm{Rtol}_r\geq0$; hence
$\mathrm{sc}_r>0$ for finite states even when a component passes through zero.
For a component of magnitude $M_r$, the scale is
$\mathrm{Atol}_r+\mathrm{Rtol}_r M_r$: the absolute term supplies a floor near
zero, and the relative term increases the permitted error with the magnitude.

HNW's Eqs.~(4.10)--(4.11) motivate componentwise scales and then combine them
using a root-mean-square measure:
\begin{equation}\label{Eq:AdaptErrorNorm}
 \mathrm{err}
 =\left[\frac1d\sum_{r=1}^d
        \left(\frac{\delta^r}{\mathrm{sc}_r}\right)^2\right]^{1/2}
 =\left(\frac{z_1^2+\cdots+z_d^2}{d}\right)^{1/2}.
\end{equation}
The explicit finite sum here defines the norm across components; it is not
an unscaled norm of quantities with incompatible units.

For example, take two components with
\[
\begin{aligned}
 \max(|y_n^1|,|(y_{n+1}^{\mathrm{high}})^1|)&=10,
 &\mathrm{Atol}_1&=10^{-6},&\mathrm{Rtol}_1&=10^{-3},\\
 \max(|y_n^2|,|(y_{n+1}^{\mathrm{high}})^2|)&=0,
 &\mathrm{Atol}_2&=10^{-6},&\mathrm{Rtol}_2&=10^{-3}.
\end{aligned}
\]
Then
\[
\begin{aligned}
 \mathrm{sc}_1&=10^{-6}+10^{-3}(10)=0.010001,\\
 \mathrm{sc}_2&=10^{-6}+10^{-3}(0)=10^{-6}.
\end{aligned}
\]
If $\delta^1=0.0050005$ and $\delta^2=5\times10^{-7}$, the complete reduction is
\[
 z_1=\frac{0.0050005}{0.010001}=\frac12,\qquad
 z_2=\frac{5\times10^{-7}}{10^{-6}}=\frac12,
\]
\[
 \mathrm{err}=\sqrt{\frac{(1/2)^2+(1/2)^2}{2}}=\frac12.
\]
This attempted step passes the test $\mathrm{err}\leq1$.
RMS acceptance does \emph{not} require every $|z_r|\leq1$:
\[
 (z_1,z_2)=(1.2,0)
 \quad\Longrightarrow\quad
 \mathrm{err}=\sqrt{\frac{1.44}{2}}\approx0.849<1.
\]
A strict componentwise test would instead use $\max_r|z_r|\leq1$.
Equation~\ref{Eq:AdaptErrorNorm} defines the RMS test used here.

\subsubsection{Deriving the next step size from the local error power}
At fixed starting data, suppose the leading error coefficient is nonzero and
the scales vary little between trial step sizes. Equations~\ref{Eq:AdaptEstimatedOrder}
and \ref{Eq:AdaptErrorNorm} give the local model
\[
 \mathrm{err}(h)\approx K h^{q+1},\qquad K>0.
\]
The scale changes with the proposed endpoint, so this is a prediction used to
choose a trial step, not an exact law. To target unit scaled error, solve
\begin{equation}\label{Eq:AdaptOptimalStep}
\begin{aligned}
 \mathrm{err}&\approx K h^{q+1},
 &1&\approx K h_{\mathrm{opt}}^{q+1},\\
 \frac1{\mathrm{err}}&\approx
       \frac{K h_{\mathrm{opt}}^{q+1}}{K h^{q+1}}
       =\left(\frac{h_{\mathrm{opt}}}{h}\right)^{q+1},\\
 h_{\mathrm{opt}}&=h\,\mathrm{err}^{-1/(q+1)}.
\end{aligned}
\end{equation}
This is HNW Eq.~(4.12). DOPRI5 has $q=4$, so the basic exponent is
$1/(q+1)=1/5$, not $1/6$: the controlling difference starts at $h^5$.

Introduce a safety factor $0<\eta<1$ and growth bounds
$0<f_{\min}<1<f_{\max}$. HNW Eq.~(4.13) is
\begin{equation}\label{Eq:AdaptBoundedStep}
 \rho=\min\!\left(f_{\max},
          \max\!\left(f_{\min},\eta\,\mathrm{err}^{-1/(q+1)}\right)\right),
 \qquad h_{\mathrm{new}}=h\rho.
\end{equation}
Before clipping, substitution into the local model shows why the safety factor
leaves room below the acceptance boundary:
\[
 K h_{\mathrm{new}}^{q+1}
 =K h^{q+1}\eta^{q+1}\mathrm{err}^{-1}
 \approx\mathrm{err}\,\eta^{q+1}\mathrm{err}^{-1}
 =\eta^{q+1}<1.
\]
For $\mathrm{err}=0$, do not evaluate a negative power: choose the largest
permitted increase. A zero measured difference can also arise from roundoff
or a vanishing leading coefficient; it does not prove zero true error.

\paragraph{A rejected step, calculated explicitly.}
For the basic DOPRI5 controller with memory disabled, take
$q=4$, $\eta=0.9$, $h=0.1$, and $\mathrm{err}=32$. Then
\[
\begin{aligned}
 h_{\mathrm{opt}}&=0.1(32)^{-1/5}=0.1/2=0.05,\\
 h_{\mathrm{new}}&=0.9(0.05)=0.045,\\
 \mathrm{err}_{\mathrm{predicted}}
 &=32\left(\frac{0.045}{0.1}\right)^5
   =32(0.45)^5=0.9^5=0.59049.
\end{aligned}
\]
The original step is rejected because $32>1$. Recompute the stages from the
same $(x_n,\mathbf y_n)$ using $0.045$; accept only if the newly computed
error passes the test. The prediction $0.59049$ is not a substitute for that
new calculation.

\subsubsection{Acceptance changes the state; rejection changes only the trial step}
Check both finiteness and the scalar error. For a finite trial,
\begin{equation}\label{Eq:AdaptAcceptReject}
\begin{array}{lll}
 \mathrm{err}\leq1:
 &x_{n+1}=x_n+h,
 &\mathbf y_{n+1}=\mathbf y_{n+1}^{\mathrm{high}},\\[2pt]
 \mathrm{err}>1:
 &x_n\ \text{unchanged},
 &\mathbf y_n\ \text{unchanged; retry with }h_{\mathrm{new}}.
\end{array}
\end{equation}
If the trial state or error is nonfinite, reject the trial and use $h_{\mathrm{new}}=f_{\min}h$. After a rejected attempt, the first accepted
retry is not allowed to propose immediate step growth; its proposed next
magnitude is capped by the accepted retry's magnitude.

The accepted state is already the higher-order output. In particular,
\[
 \mathbf y_{n+1}=\mathbf y_{n+1}^{\mathrm{high}}
 \quad\text{does not mean}\quad
 \mathbf y_{n+1}=\mathbf y_{n+1}^{\mathrm{high}}+\boldsymbol\delta.
\]
The estimate controls the trial step; it is not an extra correction added to
the fifth-order DOPRI5 state. Nor does $\mathrm{err}\leq1$ bound the accumulated
global error since the original initial time.

\subsubsection{PI step-size control and its memory}
The plain rule in Eq.~\ref{Eq:AdaptBoundedStep} uses only the current error.
A controller with memory also uses the previous accepted step's scaled error
$\mathrm{err}_{\mathrm{old}}>0$:
\begin{equation}\label{Eq:AdaptPIProposal}
 \rho_* =\eta\,\mathrm{err}^{-\alpha}
                   \mathrm{err}_{\mathrm{old}}^{\beta},
 \qquad
 \log\rho_* =\log\eta-\alpha\log\mathrm{err}
                         +\beta\log\mathrm{err}_{\mathrm{old}}.
\end{equation}
The logarithmic expression shows separately the current-error response and
the remembered-error contribution. This is a controller modification, not
another cancellation of the local truncation error. Setting $\beta=0$ and
$\alpha=1/(q+1)$ recovers the basic rule. See HNW I, printed p.~168, for the
reference to PI control in Volume~II \cite{HNW1993}.

One choice of controller parameters for DOPRI5 is
\begin{equation}\label{Eq:AdaptDopriPIConstants}
\begin{aligned}
 \eta&=0.9,& f_{\min}&=0.2,& f_{\max}&=10,\\
 \beta&=0.04,&
 \alpha&=\frac15-0.75\beta
        =0.2-0.03=0.17.
\end{aligned}
\end{equation}
Thus the accepted-step proposal before clipping is
$h_{\mathrm{new}}=0.9h\,\mathrm{err}^{-0.17}
\mathrm{err}_{\mathrm{old}}^{0.04}$.
These are adjustable controller parameters; they are not part of the
Runge--Kutta order conditions.

The same step-size rule can be written using a reciprocal factor. For an
accepted trial,
\begin{equation}\label{Eq:AdaptPIReciprocal}
\begin{aligned}
 F_*&=\frac{\mathrm{err}^{\alpha}}
           {\eta\,\mathrm{err}_{\mathrm{old}}^{\beta}}
       =\frac1{\rho_*},\\
 F&=\max\!\left(\frac1{f_{\max}},
           \min\!\left(\frac1{f_{\min}},F_*\right)\right),\\
 h_{\mathrm{new}}&=\frac hF.
\end{aligned}
\end{equation}
For positive finite errors this is algebraically the same clipping as
$h\min(f_{\max},\max(f_{\min},\rho_*))$. The division order can change the
last floating-point bits. If the trial was preceded by a rejection, also
cap $|h_{\mathrm{new}}|\leq|h|$.

For a rejected finite trial, drop the memory term:
\begin{equation}\label{Eq:AdaptPIRejected}
 F_{\mathrm{rej}}=\min\!\left(\frac1{f_{\min}},
                         \frac{\mathrm{err}^{\alpha}}{\eta}\right),
 \qquad h_{\mathrm{new}}=\frac h{F_{\mathrm{rej}}}.
\end{equation}
For the DOPRI5 parameter choice in Eq.~\ref{Eq:AdaptDopriPIConstants},
$\alpha=0.17$ is retained even on rejection; the
$1/5$ example above described the memory-disabled choice $\beta=0$.
Only acceptance updates the remembered error:
\[
 \mathrm{err}_{\mathrm{old}}\leftarrow\max(\mathrm{err},10^{-4}).
\]
The initial remembered value is $10^{-4}$. This floor keeps the memory factor
away from zero. When $\mathrm{err}=0$, the reciprocal calculation for these
positive-$\alpha$ choices selects $f_{\max}$ without division by zero in
$h/F$.

\subsubsection{Choosing a first step before an error estimate exists}
Hairer I, printed p.~169, Eq.~(4.14), proposes an Euler probe to estimate the
initial rate of change. The following
formulas describe its positive step magnitudes; a backward integration gives
both the probe and the final step the integration direction's sign.

At $(x_0,\mathbf y_0)$, compute $\mathbf f_0=\mathbf f(x_0,\mathbf y_0)$ and
freeze the initial scales
\[
 \mathrm{sc}_{0r}=\mathrm{Atol}_r+\mathrm{Rtol}_r|y_0^r|.
\]
For this starting-step calculation, define
\begin{equation}\label{Eq:AdaptInitialNorms}
 D_0=\left[\sum_r\left(\frac{y_0^r}{\mathrm{sc}_{0r}}\right)^2\right]^{1/2},
 \qquad
 D_1=\left[\sum_r\left(\frac{f_0^r}{\mathrm{sc}_{0r}}\right)^2\right]^{1/2}.
\end{equation}
These norms omit the RMS divisor $d$. For the same scales and probe,
write the book's RMS quantities as $d_j^{\mathrm{RMS}}$. Then
\[
 D_j=\sqrt d\,d_j^{\mathrm{RMS}}\quad(j=0,1,2),\qquad
 \frac{D_0}{D_1}=\frac{d_0^{\mathrm{RMS}}}{d_1^{\mathrm{RMS}}}.
\]
The divisor cancels in the first-guess ratio, but not in every later formula.
For $D_0,D_1>10^{-5}$, the first guess satisfies
\[
 h_0=0.01\frac{D_0}{D_1}
 \quad\Longrightarrow\quad
 \left\|\left(\frac{h_0f_0^r}{\mathrm{sc}_{0r}}\right)_r\right\|_2
 =h_0D_1=0.01D_0.
\]
Thus the scaled Euler increment is one percent of the scaled starting state.
If either norm is at most $10^{-5}$, use $h_0=10^{-6}$ instead. Cap $h_0$
by the prescribed maximum step before the probe.

Make one Euler prediction and evaluate its derivative:
\begin{equation}\label{Eq:AdaptEulerProbe}
 \mathbf y_{\mathrm E}=\mathbf y_0+h_0\mathbf f_0,
 \qquad
 \mathbf f_{\mathrm E}=\mathbf f(x_0+h_0,\mathbf y_{\mathrm E}).
\end{equation}
A first-order expansion displays what this probe measures:
\begin{align*}
 \frac{f_{\mathrm E}^r-f_0^r}{h_0}
 &=\frac{\partial f^r}{\partial x}(x_0,\mathbf y_0)
   +\frac{\partial f^r}{\partial y^j}(x_0,\mathbf y_0)f_0^j+O(h_0)\\
 &=\frac{d^2y^r}{dx^2}(x_0)+O(h_0).
\end{align*}
Set
\[
 D_2=\frac1{h_0}\left[\sum_r
       \left(\frac{f_{\mathrm E}^r-f_0^r}{\mathrm{sc}_{0r}}\right)^2
       \right]^{1/2},\qquad M=\max(D_1,D_2).
\]
This uses one extra right-hand-side evaluation beyond $\mathbf f_0$.
For $M>10^{-15}$ the second guess solves
\begin{equation}\label{Eq:AdaptSecondGuess}
 h_1^kM=0.01
 \quad\Longrightarrow\quad
 h_1=\left(\frac{0.01}{M}\right)^{1/k}.
\end{equation}
For a fixed exponent $k$, outside the fallback and clipping branches,
the norm change alone gives
\[
\begin{aligned}
 h_{1,\mathrm{sum}}
 &=\left(\frac{0.01}{\sqrt d\,M_{\mathrm{RMS}}}\right)^{1/k}\\
 &=d^{-1/(2k)}h_{1,\mathrm{RMS}}.
\end{aligned}
\]
The fallback thresholds may also select different branches. Separately,
the book writes $k=p+1$ in this heuristic. Alternative choices are $k=5$
for DOPRI5 and $k=8$ for DOP853, matching the error powers used for step-size
control. If $M\leq10^{-15}$, use
$h_1=\max(10^{-6},10^{-3}h_0)$ instead. The proposed initial magnitude is
\begin{equation}\label{Eq:AdaptInitialStep}
 h_{\mathrm{start}}=\min(100h_0,h_1,h_{\max}).
\end{equation}
These starting guesses are empirical numerical rules in the chosen independent
variable's units; $D_1$ and $D_2$ involve derivatives of different orders.
They are not an exact truncation-error identity or invariant under arbitrary
rescaling of the independent variable. The resulting RK trial must still pass
the ordinary error test.

For a scalar check, let $y'=-y$, $y_0=1$, $\mathrm{Atol}=10^{-8}$,
$\mathrm{Rtol}=0$, and let the maximum step be large enough not to clip.
Then
\[
\begin{aligned}
 D_0&=D_1=10^8, &h_0&=0.01,\\
 y_{\mathrm E}&=1-0.01=0.99,
 &f_{\mathrm E}-f_0&=-0.99-(-1)=0.01,\\
 D_2&=\frac{0.01/10^{-8}}{0.01}=10^8,
 &M&=10^8.
\end{aligned}
\]
For DOPRI5, $k=5$, so
\[
 h_1=\left(\frac{10^{-2}}{10^8}\right)^{1/5}
     =(10^{-10})^{1/5}=10^{-2},\qquad
 h_{\mathrm{start}}=\min(1,0.01)=0.01.
\]

\subsubsection{Where DOP853 changes the error calculation}
DOP853 retains an eighth-order state, but its controller does not simply use
an eighth-minus-seventh-order difference. It combines fifth- and third-order
embedded differences. Define
\[
 \boldsymbol\delta_5=\mathbf y^{(8)}-\mathbf y^{(5)},\qquad
 \boldsymbol\delta_3=\mathbf y^{(8)}-\mathbf y^{(3)}.
\]
These use retained minus embedded, the opposite sign to HNW Eq.~(10.16).
The sign does not affect the squared sums. Define
\[
 S_5=\sum_r\left(\frac{\delta_5^r}{h\,\mathrm{sc}_r}\right)^2,
 \qquad
 S_3=\sum_r\left(\frac{\delta_3^r}{h\,\mathrm{sc}_r}\right)^2.
\]
Combine these scaled differences as
\begin{equation}\label{Eq:AdaptDop853Composite}
 \mathrm{err}_{853}
 =\frac{|h|S_5}{\sqrt{d(S_5+0.01S_3)}}.
\end{equation}
For finite differences, if $S_5=S_3=0$, take $\mathrm{err}_{853}=0$.
Generically, at fixed scales with nonzero leading coefficients,
\[
 \boldsymbol\delta_5=O(h^6),\quad
 \boldsymbol\delta_3=O(h^4)
 \quad\Longrightarrow\quad
 S_5\sim A h^{10},\quad S_3\sim B h^6,\qquad A,B>0.
\]
Substitution into Eq.~\ref{Eq:AdaptDop853Composite} gives
\[
 \mathrm{err}_{853}
 \sim\frac{A h^{11}}{\sqrt{0.01dB h^6}}
 =\frac{A}{\sqrt{0.01dB}}h^8.
\]
This explains the controller's exponent $1/8$; the accepted eighth-order
solution itself has local error $O(h^9)$. If the third-order difference is
unusually small, the dominant balance changes and the generic $h^8$ argument
need not apply. One controller parameter choice for DOP853 is $\beta=0$, $\alpha=1/8$,
$\eta=0.9$, $f_{\min}=0.333$, and $f_{\max}=6$.
The eighth-order method is constructed in Hairer I, Section~II.5, printed
pp.~181--185. Its fifth-/third-order error estimator is developed in
Section~II.10, printed pp.~254--255, Eqs.~(10.16)--(10.17) \cite{HNW1993}.
Equation~\ref{Eq:AdaptDop853Composite} gives the scaled vector form of this
composite estimate.

\subsubsection{What an adaptive step actually does}
For DOPRI5 the calculation proceeds as
\[
 \{\mathbf k_i\}_{i=1}^7
 \ \longrightarrow\
 \left(\mathbf y_{n+1}^{\mathrm{high}},
       \boldsymbol\delta=h\sum_i\varepsilon_i\mathbf k_i\right)
 \ \longrightarrow\ \mathrm{err}.
\]
The error norm compares this attempt with the chosen tolerance scales.
Acceptance advances to the higher-order endpoint; rejection leaves the
starting state fixed and recalculates with a smaller trial step. The controller
uses the measured error to propose the next step size, and every new proposal
is tested again. This feedback calculation is what the adaptive solver adds
to a fixed-step RK update.
